\documentclass{resume} % Use the custom resume.cls style
\usepackage{enumitem} % for item spacing
\usepackage{scalerel} % for inline images

\usepackage[left=0.4 in,top=0.2in,right=0.4 in,bottom=0.2in]{geometry} % Document margins
\usepackage{xcolor}
\newcommand{\tab}[1]{\hspace{.2667\textwidth}\rlap{#1}} 
\newcommand{\itab}[1]{\hspace{0em}\rlap{#1}}
\name{JOHN GIORSHEV} % Your name

\address{
    \scalerel*{\includegraphics{images/location.png}}{|} \href{https://www.google.com/maps/place/Greater+Toronto+Area}{Toronto, ON}\space\space\space\space\space\space
    \scalerel*{\includegraphics{images/email.png}}{|} \href{mailto:johngiorshev1@gmail.com}{johngiorshev1@gmail.com}\space\space\space\space\space
    \scalerel*{\includegraphics{images/phone.png}}{|} \href{tel:+17058164454}{705-816-4454}\space\space\space\space\space\space
    \scalerel*{\includegraphics{images/linkedin.png}}{|} \href{https://www.linkedin.com/in/johngiorshev/}{johngiorshev}
}  % Your phone number, email, linkedin, and (optional) website

\begin{document}

Legally authorized to work in the United States.

\begin{rSection}{EMPLOYMENT EXPERIENCE}
    \begin{itemize}
        \item {\bf ECE1724H TA (Embedded development)} - University of Toronto \hfill {Sep 2023 - Jan 2024}
        \begin{itemize}[topsep=-10pt]
            \setlength\itemsep{-0.5em}
                \item[\textbullet] {\bf Special Topics in Computing} - \textcolor{darkgray}{Wearable AI, Building Superintelligence for People}
            \end{itemize}
        \item {\bf Software Developer Intern} - SpaceRyde \hfill {Sep 2021 - Sep 2022}
        \begin{itemize}[topsep=-10pt]
            \setlength\itemsep{-0.5em}
            \item[\textbullet] {\bf Embedded Development} - \textcolor{darkgray}{ESP32 embedded development using PlatfromIO. MQTT.}
            \item[\textbullet] {\bf User Interfaces} - \textcolor{darkgray}{Wrapped ROS2 components via QT GUIs and FSMs for human interface and control.}
            \item[\textbullet] {\bf Telemetry} - \textcolor{darkgray}{Data acquisition via LabJack. Prioritization, packetization, and streaming of sensors and metrics. Downstream dashboards with TimescaleDB and Grafana.}
            \item[\textbullet] {\bf Simulator} - \textcolor{darkgray}{Decreased computation time of a 6DOF rigid body simulator \textbf{by 50x} using Eigen lib and parallelism.
            Integrated with ROS2 and GeographicLib.}
        \end{itemize}
        \item {\bf Cyber Threat Intelligence Engineering Co-op} - Bell Canada Enterprises \hfill {Summer 2020}
        \begin{itemize}[topsep=-10pt]
            \setlength\itemsep{-0.5em}
            \item[\textbullet] \textcolor{darkgray}{Integrated a MISP server with an Elasticsearch database using RESTful APIs. Elastic Stack. Kafka.}
            \item[\textbullet] \textcolor{darkgray}{Improved Bell's security information infrastructure. Collaborated using Confluence and Jira.}
            \item[\textbullet] \textcolor{darkgray}{Mattermost (slack) notification bot allowing faster team response to latest disseminated exploits.}
        \end{itemize}
    \end{itemize}
\end{rSection}

\begin{rSection}{Projects}
    \begin{itemize}
        \setlength\itemsep{-0.2em}
        \item {\bf Choose} - \href{https://github.com/jagprog5/choose/}{jagprog5/choose} \hfill C++\\
        \textcolor{darkgray}{Created a stream manipulator that can be \textbf{up to 100x faster} than standard tools including GNU sed.
        Uses the PCRE2 regex C API for pattern matching and Ncurses C API for TUIs.}
        \item {\bf Key Value} - \href{http://www.keyvalue.ca/}{keyvalue.ca} (hosting OpenVPN access server at \href{https://www.keyvalue.ca:943}{keyvalue.ca:943})\hfill Rust\\
        \textcolor{darkgray}{HTTPS enabled self-hosted web site for users to store key value pairs. Docker Compose. Actix Web HTTP backend integrated with SQLite. TLS termination with NGINX and Let's Encrypt for secure ingress.}
        \item {\bf Game Engine} - \href{https://github.com/jagprog5/game_engine/}{jagprog5/game\_engine} \hfill Rust\\
        \textcolor{darkgray}{(in progress) Creating a 2D game engine built on SDL2 featuring full persistance (serde).}
        \item {\bf Sparse Distributed Representation} - \href{https://github.com/jagprog5/SDR/}{jagprog5/sdr} \hfill C++\\
        \textcolor{darkgray}{Sparse array library. Template metaprogramming and specializations for data structures. Profile guided optimization. Integrated with CMake, CI/CD: static analyzers, fuzzy testing for complete code coverage.}
        \item {\bf AWS Log Uploader} \hfill Python\\
        \textcolor{darkgray}{Created a multithreaded utility that tracks a log folder and uploads to AWS CloudWatch using boto3.\\
        Robust to sudden interruption and prolonged internet outage.}
        \item {\bf Data Analytics} - \href{https://github.com/jagprog5/resume/blob/main/data-analytics-projects.md}{jagprog5/data-analytics} \hfill Python\\
        \textcolor{darkgray}{Implemented multi-layer perceptrons using PyTorch and C. Implemented hopfield networks.\\
        Used various sklearn models on Kaggle datasets. Created a shell wrapper of Matplotlib for data visualization.}
    \end{itemize}
\end{rSection}

\begin{rSection}{Design Teams}
    \begin{itemize}
        \item {\bf QMind} - Queen's AI/ML Software Team \hfill {Sep 2019 - May 2020}\\
        \textcolor{darkgray}{Presented at \href{https://medium.com/@cameronmackinn/cucai-2020-c821c15f3ab4}{CUCAI 2020}. Used OpenCV and Flask to create a computer vision API.}
        \item {\bf Network Security Team} \hfill {Sep 2018 - May 2019}\\
        \textcolor{darkgray}{Penetration testing on virtual machines with Metasploit. Exploited apps in a sandboxed environment.}
\end{itemize}
\end{rSection}

\begin{rSection}{EDUCATION}
    \begin{itemize}
        \item {\bf Computer Engineering} - Queen's University (First Class Honours) \hfill {Sep 2018 - May 2023}
    \end{itemize}
\end{rSection}

\end{document}
